
\subsubsection{Fiber Section}
\begin{Prompt}
 If the element has warping freedoms (Type 1), the section response is 
$$
(\gamma,\kappa,\alpha,\alpha')\mapsto (n,m,w,v)
$$
with generalized warping resultants $\bm{w}$ and $\bm{v}$, but if the element is Type 3 (classical with 6 dofs) then this full section response is condensed via the trace matrix. 
This works on two ends. 
The element gives $\bm{\gamma},\bm{\kappa}$, and the middleman reconstructs $\alpha,\alpha'$ before invoking the generalized section. 
When the general section returns with the full $(\bm{n},\bm{m},\bm{w},\bm{v})$ the middleman condenses it to a 6D vector resembling n,m before passing it back up to the element. 
This can be done in one of two ways: in a Pevrov or Bubnov Galerkin style. 
We always assume these Type 3 elements enforce the classical equilibrium operator $\mathbb{B}$ and its adjoint $\mathbb{B}^*$. 
If the trace is class II (a subset of class I), we can use Bubnov-Galerkin and pass return the conjugate resultants to the trace's $\bm e$. 
If the trace is only class I, then the conjugate resultants are not exactly equilibriated by $\mathbb{B}$, and we must instead return the classical resultants $(\bm n,\bm m)$.
\end{Prompt}


The warping basis is, in terms of the functions \eqref{eq:def-eight-shapes}:
\[
\bm{\Phi}=\Big[
\ \bm{W}_{(\varepsilon)},
\ \bm{W}_{(a)},
\ \bm{\Pi}_{(b)},
\ \FrameBasis\,\varphi,
\ \FrameBasis\otimes\bm{\psi}
\Big].
\]
\begin{Aside}
$$
\begin{aligned}
\bm{\Phi} = 
\left[\begin{array}{c|cc|cc|c|cc} %|cc|c}
-\nu \bm{r} & \bm{W}_{(a)} & &\bm{\Pi}_{(b)} & & \FrameBasis\WarpTwistIesan & \FrameBasis\otimes\bm{\WarpShearIesan}  \\ \hline 
%\FrameBasis\otimes\bm{r} &  \FrameBasis \\ \hline
~~~~ 0 & 0 & 0 & 0 & 0 &  \varphi & \psi_{y} & \psi_{z} \\ %&  y & z & 1 \\
-\nu y & - & - & - & - &  0 & 0 & 0                     \\ %&  0 & 0 & 0 \\
-\nu z & - & - & - & - &  0 & 0 & 0                     \\ %&  0 & 0 & 0 \\
\end{array}\right]
\end{aligned}
$$
\end{Aside}

\paragraph{Amplitudes}

\begin{equation}
\begin{aligned}
\alpha_{\WarpTwistIesan}&=a_\vartheta+c_\vartheta&&=\tilde{\Twist} = (\ShiftB^{\top}\CenterShearLocation + \ShiftT)\cdot\bm{\gamma} + (\CenterShearLocation\cdot\ShiftK + \ShiftScaleT) \Twist \\
\bm{\alpha}_{\WarpShearIesan}&=\bm{b}_{\Section}&&=\ShiftB\bm{\gamma} + \ShiftK\Twist \\
\bm{\alpha}_{r}&=-\GaugeShearA  &&=-\bm{\gamma}\\.
\end{aligned}
\label{eq:identify-alpha-trace}
\end{equation}
%


\subsubsection{Implementation}

\paragraph{Definition.}
At the time a model is defined the section is supplied with the following data:
\begin{description}
    \item[Fiber data.] At each integration point $i$ with position $\bm{r}_i$, 
    provide the values of the particular warping functions \(\WarpTwistIesan\) and \(\bm{\WarpShearIesan}\).
    
    \item[Global data.] The inverse trace matrix $\TraceDpDeRoot$ (\cref{tab:T0-inv-blocks}).
\end{description}

\paragraph{Evaluation.} 
Given beam strains $\bm{e} = (\bm{\gamma}, \bm{\kappa})$:
\begin{enumerate}
\item At each fiber $i$, evaluate the 3D strains directly:
\RepeatEquation{eq:eps-B-e}
\end{enumerate}


\begin{Old}

\textbf{Pre-processing.}
For an implementable trace, the section stores:
\begin{description}
    \item[Fiber Data.] At each integration point $i$ with position $\bm{r}_i$:
    \begin{itemize}
    \item Torsion warping gradient: $\nabla\WarpTwistIesan(\bm{r}_i)$
    \item Flexural shear tensor: 
    $\mathbf{B}_{(b)} \triangleq \PoissonFlexure(\bm{r}_i) + (\nabla\bm{\WarpShearIesan}(\bm{r}_i))^\top$
    \end{itemize}

    \item[Global data.] The conversion matrix 
$\mathbf{X}^{-1} = \mathbf{L}_0\,\TraceDpDeRoot$, which maps beam strains 
to local parameters.
\end{description}
\end{Old}